\section{Method}
\label{sec:method}

\subsection{Generalizability Theory}
\label{sec:gtheory}

Generalizability theory (G-theory) treats each measurement as one observation from a universe of admissible conditions defined by the crossed facets of an evaluation design \cite{shavelson1991generalizability, brennan2001generalizability}. The \emph{object of measurement} is the entity whose consistency is of interest. In our setting, test items ($p$) serve as the object of measurement, and the goal is to determine how reliably item-level scores generalize across the methodological conditions under which they are collected.

A \emph{G-study} decomposes total observed variance into components $\sigma^2_x$ for each facet $x$ and their interactions. Each component quantifies how much variability that source contributes to the overall measurement. Large interaction components (e.g., $\sigma^2_{pM}$ for item-by-model) indicate that facets do not affect items uniformly but instead create differential difficulty profiles. We estimate variance components via Henderson Method~I, an unbiased ANOVA-based estimator for balanced designs \cite[Ch.~3]{brennan2001generalizability}, with 1,000-resample bootstrap confidence intervals \cite{li2023bootstrap}.

A \emph{D-study} projects the generalizability coefficient under hypothetical designs with different numbers of conditions per facet:
\begin{equation}
G_\text{item} = \frac{\sigma^2_p}{\sigma^2_p + \sigma^2_\text{rel}}
\label{eq:gcoeff}
\end{equation}
where $\sigma^2_\text{rel}$ aggregates all item-by-facet interaction variances, each divided by the number of conditions for the corresponding facets \cite[Ch.~4]{brennan2001generalizability}. This prospective capability distinguishes G-theory from classical ANOVA: rather than only attributing variance post hoc, a D-study enables optimization of evaluation designs \emph{before} data collection \cite{shavelson1989gtheory}. Throughout this paper, we use the following notation: $p$ = item, $M$ = model, $\pi$ = numerical precision, $\delta$ = prompt template, $s$ = random seed, $\tau$ = sampling temperature, $o$ = answer ordering, and $\sigma^2_x$ = variance component for source $x$.

\subsection{Experimental Design: MMLU}
\label{sec:exp001}

We design a fully crossed factorial experiment (exp-001) that simultaneously varies seven facets of LLM evaluation on MMLU. Table~\ref{tab:design_matrix} summarizes the complete design matrix, and Figure~\ref{fig:design_diagram} illustrates the crossing structure.

\paragraph{Facets.} Items ($p$): 200 MMLU questions stratified across 10 subjects (20 per subject). The stratified sampling ensures that variance estimates are not dominated by a single knowledge domain; the 10 subjects span STEM, humanities, social sciences, and professional disciplines. Models ($M$): Llama-3.1-8B-Instruct, Gemma-2-9B-it, Mistral-7B-Instruct-v0.3, and Qwen3-8B, a purposive sample from four major architecture families treated as a quasi-random facet (\S3.6). Numerical precision ($\pi$): BF16, FP16, and FP32, all evaluated with \texttt{enforce\_eager=True} for fair comparison across precision levels. Prompt template ($\delta$): six semantically equivalent paraphrases of the evaluation instruction \cite{sclar2024formatting}. Random seed ($s$): six fixed values (42, 123, 456, 789, 1024, 2048). Sampling temperature ($\tau$): three levels (0.0, 0.3, 0.7) spanning deterministic to moderately stochastic decoding. Answer ordering ($o$): four cyclic rotations of multiple-choice option positions \cite{pezeshkpour2024ordering}.

\paragraph{Sample sizes.} The cross-model BF16 balanced subset yields 345,600 records (4 models $\times$ 200 items $\times$ 6 prompts $\times$ 6 seeds $\times$ 3 temperatures $\times$ 4 orderings). A single-model analysis (Llama) spans all three precisions for 239,200 records (1,196 conditions $\times$ 200 items). The BF16 subset serves as the primary cross-model sample because BF16 is the standard inference precision for these architectures. Gemma-2 does not support FP16 due to vLLM numerical instability guards; cross-model analyses therefore use BF16 to maintain a balanced design.

\begin{table}[t]
\centering
\caption{Fully crossed G-theory design. Exp-001 evaluates MMLU across all facet levels; exp-002 extends to four benchmarks with reduced crossing. $^\dagger$BF16 balanced subset for cross-model analysis.}
\label{tab:design_matrix}
\small
\begin{tabular}{llrr}
\toprule
Facet & Levels & exp-001 & exp-002 \\
\midrule
Item ($p$) & Stratified sample & 200 & 200$\times$4 \\
Model ($M$) & 4 architecture families & 4 & 4 \\
Precision ($\pi$) & BF16 / FP16 / FP32 & 3 & 1 \\
Prompt ($\delta$) & Paraphrases & 6 & 6 \\
Seed ($s$) & Fixed set & 6 & 6 \\
Temperature ($\tau$) & 0.0 / 0.3 / 0.7 & 3 & 2 \\
Ordering ($o$) & Cyclic rotations & 4 & 4 \\
\midrule
\multicolumn{2}{l}{Total records} & 345,600$^\dagger$ & 921,600 \\
\bottomrule
\end{tabular}
\end{table}

\begin{figure}[t]
\centering
\includegraphics[width=\columnwidth]{fig2_design_diagram}
\caption{Fully crossed G-theory design for LLM evaluation. Each observation is a binary correctness score determined by the crossing of seven facets: item ($p$), model ($M$), precision ($\pi$), prompt ($\delta$), seed ($s$), temperature ($\tau$), and ordering ($o$).}
\label{fig:design_diagram}
\end{figure}

\subsection{Cross-Benchmark Extension}
\label{sec:exp002}

To test whether the variance structure observed on MMLU generalizes across evaluation formats, exp-002 extends the design to four benchmarks: MMLU, ARC-Challenge, and HellaSwag (multiple-choice), and GSM8K (free-form generation). The design crosses 4 models $\times$ BF16 $\times$ 2 temperatures (0.0, 0.7) $\times$ 6 prompts $\times$ 6 seeds $\times$ 4 orderings $\times$ 200 items per benchmark, totaling 921,600 records. This reduced crossing fixes precision at BF16 and uses two temperature levels to concentrate statistical power on format-conditional differences. For GSM8K, the ordering facet permutes few-shot exemplar sequences rather than answer options. The two temperature levels bracket the most common evaluation configurations: deterministic ($\tau = 0.0$) and a moderate stochastic setting ($\tau = 0.7$) representative of practical deployment.

% [PLACEHOLDER: 8 models × 5 benchmarks when expansion complete]

\subsection{Evaluation Protocol}
\label{sec:protocol}

All evaluations use generation-based scoring: models produce free-text responses from which correctness is determined by answer extraction, rather than by ranking answer options via token log-likelihood. This design choice follows \citet{hua2025flaw}, who show that log-likelihood scoring conflates model confidence with surface-level token biases. We report two metrics: binary correctness (primary, 1/0 per item) and text-exact-match (secondary). Binary correctness accepts any response containing the correct answer letter or numerical value after standard normalization; text-exact-match requires the generated string to match the reference answer character-for-character.

Inference runs on vLLM~0.9.2 with default settings and without batch-invariant kernels. This configuration reflects standard practitioner conditions and preserves CUDA-level nondeterminism from floating-point non-associativity as a measurable variance source \cite{messina2026background, thinkingmachines2025batch}. Few-shot prompting follows standard protocols: 5-shot for MMLU and ARC-Challenge, 8-shot for GSM8K, and 10-shot for HellaSwag.

\subsection{Binary Data Assumptions}
\label{sec:binary}

Standard G-theory assumes continuous, approximately normal outcome data \cite{shavelson1991generalizability}. Our binary correctness scores (0/1) violate this assumption. We establish robustness through four complementary analyses.

\paragraph{Simulation.} We generate continuous latent scores from a probit model matched to our 1,296-condition design, then dichotomize at the empirical accuracy threshold. The binary decomposition yields a generalizability coefficient of 0.732, closely approximating the continuous-data coefficient of 0.741 ($\Delta = 0.009$). Henderson Method~I identifies \texttt{item\_id} as the dominant variance source in all 500 simulation runs.

\paragraph{GLMM complement.} A generalized linear mixed model (GLMM) with marginal logit link produces variance component rankings that perfectly preserve all seven grouped component orderings ($\rho = 1.0$) relative to Henderson Method~I \cite{jiang2018bayesian, tuerlinckx2006glmm}.

\paragraph{Accuracy-range stability.} The \texttt{item\_id} variance share remains stable across subsets with mean accuracy from 0.62 to 0.72 ($r = -0.048$). This stability rules out distortion from binary floor or ceiling effects at observed accuracy levels.

\paragraph{Estimator guarantee.} Henderson Method~I produces unbiased variance component estimates for balanced designs regardless of the outcome distribution \cite[Ch.~12]{brennan2001generalizability}. Our fully crossed design satisfies this balance condition.

\subsection{Model-as-Facet Justification}
\label{sec:model_facet}

Treating four purposively selected models as a quasi-random facet requires empirical validation. We provide three forms of evidence.

D-study cross-validation: variance components estimated from all $\binom{4}{3} = 4$ three-model subsets predict the four-model generalizability coefficient. The predicted $G = 0.816$ for $n_\text{model} = 4$ closely matches the observed $G = 0.820$ ($\Delta = 0.004$). This agreement confirms that the model facet behaves as expected under random-effects assumptions.

Leave-one-model-out (LOMO) stability: all four three-model subsets preserve \texttt{item\_id} as the dominant component and rank \texttt{model$\times$item} as the second-largest ($G_\text{item}$ range: 0.747--0.793). The consistency of rankings across all subsets indicates that no single model disproportionately drives the observed variance structure.

Forward projection: D-study estimates predict $G = 0.893$ at eight models. This result supports the quasi-random facet assumption for the broader population of instruction-tuned 7--9B architectures.

% [PLACEHOLDER: 8-model expansion strengthens to ~0.893 observed]
